\clearpage
\SourceLeaf{66}
\pdfbookmark[1]{Disquisitio de elementis ellipticis Palladis — Anzeige 1810}{anzeige-palladis1810}
\begin{center}\rule{.75\textwidth}{.3pt}\\
Göttingische gelehrte Anzeigen. 13. Dezember 1810.\\
\rule{.75\textwidth}{.3pt}\end{center}
Am 25. November legte Professor \textsc{Gauß} der Königlichen Gesellschaft der Wissenschaften eine Abhandlung vor:
\begin{center}
\textit{Disquisitio de elementis ellipticis Palladis ex oppositionibus annorum}\\
1803, 1804, 1805, 1807, 1808, 1809.
\end{center}
Seit drei Jahren hatte der Verfasser keine neuen Rechnungen zur Bahn der Pallas angestellt. Wegen der großen Lichtschwäche des Planeten waren die Beobachtungen von 1808 sehr spärlich und mangelhaft gewesen, und zum Teil wurden sie erst später bekannt. Deshalb schien es nicht der Mühe wert, die Elemente schon damals danach zu verbessern. Erst nachdem Professor \textsc{Gauß} die Beobachtungen der Pallas von 1809 erhalten hatte, die \textsc{Bouvard} an der Kaiserlichen Sternwarte in Paris angestellt hatte, berechnete er zusammen mit der Opposition von 1809 auch die von 1808 (siehe diese Gelehrten Anzeigen vom 24. Februar 1810). Die weitere Untersuchung aller sechs bisher beobachteten Oppositionen ergab, dass eine elliptische Bahn nicht mehr ausreicht, um sie alle genau darzustellen. Das ist eine Folge der großen Störungen, die dieser Planet durch die anderen, besonders durch Jupiter, erfährt. Professor \textsc{Gauß} hat dies auf zweierlei Weise gezeigt. Zunächst berechnete er drei Systeme elliptischer Elemente, jedes aus vier Oppositionen, das erste aus denen von 1803, 1804, 1805 und 1807,
\SourceLeaf{67}%
das zweite aus denen von 1804, 1805, 1807 und 1808 und das dritte aus den Oppositionen von 1805, 1807, 1808 und 1809. Wäre die Bewegung rein elliptisch, so hätten sich zwischen diesen Systemen nur kleine Unterschiede zeigen müssen. Wir stellen die drei Systeme hier nebeneinander:
\par\nobreak
\begingroup
\newcommand{\GaussVIAngle}[6]{%
  \makebox[17pt][r]{$#1$}\makebox[4pt][l]{${}#2$}\hspace{2pt}%
  \makebox[12pt][r]{$#3$}\makebox[4pt][l]{${}#4$}\hspace{2pt}%
  \makebox[12pt][r]{$#5$}\makebox[18pt][l]{${}#6$}%
}
{\centering\small
\setlength{\tabcolsep}{8pt}
\begin{tabular}{c|l|l|l}
& \multicolumn{1}{c|}{I.} & \multicolumn{1}{c|}{II.} & \multicolumn{1}{c}{III.}\\\hline
$A$ & \GaussVIAngle{221}{^0}{39}{'}{30}{''4} & \GaussVIAngle{221}{^0}{34}{'}{56}{''7} & \GaussVIAngle{221}{^0}{23}{'}{24}{''6}\\
$B$ & \multicolumn{1}{r|}{$770.2143$} & \multicolumn{1}{r|}{$770.4467$} & \multicolumn{1}{r}{$770.9265$}\\
$C$ & \GaussVIAngle{121}{}{3}{}{11}{.4} & \GaussVIAngle{121}{}{5}{}{22}{.1} & \GaussVIAngle{120}{}{58}{}{4}{.8}\\
$D$ & \GaussVIAngle{172}{}{28}{}{56}{.9} & \GaussVIAngle{172}{}{28}{}{46}{.8} & \GaussVIAngle{172}{}{27}{}{52}{.4}\\
$E$ & \GaussVIAngle{34}{}{37}{}{41}{.0} & \GaussVIAngle{34}{}{37}{}{31}{.5} & \GaussVIAngle{34}{}{36}{}{49}{.4}\\
$F$ & $0.2450198$ & $0.2447624$ & $0.2446335$\\
$G$ & $0.4423149$ & $0.4422276$ & $0.4420473$
\end{tabular}\par}
\begin{samepage}\small
Dabei stehen die Zeichen für die folgenden Größen
\par\noindent\hspace*{.06\linewidth}%
\begin{tabular}{@{}l@{\enspace}p{.74\linewidth}@{}}
  $A$ & Epoche der mittleren Länge 1803 für den Göttinger Meridian\\
  $B$ & mittlere tägliche tropische Bewegung\\
  $C$ & Länge des Perihels 1803\\
  $D$ & Länge des aufsteigenden Knotens 1803\\
  $E$ & Neigung der Bahn\\
  $F$ & Exzentrizität\\
  $G$ & Logarithmus der großen Halbachse\\
\end{tabular}\par\smallskip
\end{samepage}
Will man eine bleibende Übereinstimmung der Erscheinungen mit der Rechnung erreichen, so wird man die Störungen der Pallas durch die anderen Planeten einrechnen müssen. Für manche Zwecke wird man aber auch künftig eine rein elliptische Bahn vorziehen, wenn sich damit die unerträglich langwierigen Rechnungen vermeiden lassen, die die große Zahl der Störungsgleichungen sonst nötig machen würde. Geht es insbesondere nur darum, den Planeten wiederzufinden, so ist es gewiss bequemer, rein elliptische Elemente auf die vier unmittelbar vorausgegangenen Oppositionen zu stützen und danach die Ephemeride zu berechnen. Nach einem Jahr kann diese noch nicht sehr weit von dem abweichen, was man am Himmel sieht. Deshalb hat Professor \textsc{Gauß} die im Oktoberheft der \textit{Monatlichen Correspondenz} abgedruckte Ephemeride für die Bewegung der Pallas bei ihrer nächsten Erscheinung nach dem obigen System III von Elementen berechnet.

\SourceLeaf{68}%
Auf eine zweite Weise hat Professor \textsc{Gauß} den Einfluss der Störungen nachgewiesen, indem er rein elliptische Elemente berechnete, die sich \textit{allen} sechs Oppositionen \textit{möglichst} genau anschließen und dennoch von den einzelnen beobachteten Orten erheblich abweichen. Wir geben hier auch dieses \textit{vierte} System von Elementen an:
\begin{center}\small
\newcommand{\GaussVIElementLeaders}{%
  .\hspace{.95em}.\hspace{.95em}.\hspace{.95em}.\hspace{.95em}.\hspace{.95em}%
  .\hspace{.95em}.\hspace{.95em}.\hspace{.95em}.\hspace{.95em}.\hspace{.95em}%
  .\hspace{.95em}.\hspace{.95em}.\hspace{.95em}.\hspace{.95em}.\hspace{.95em}.%
}
\begin{tabular}{cll}
$A$ & \GaussVIElementLeaders & \GaussVIAngle{221}{^0}{34}{'}{53}{''64}\\
$B$ & \GaussVIElementLeaders & \multicolumn{1}{r}{$770.5010$}\\
$C$ & \GaussVIElementLeaders & \GaussVIAngle{121}{}{8}{}{8}{.54}\\
$D$ & \GaussVIElementLeaders & \GaussVIAngle{172}{}{28}{}{12}{.43}\\
$E$ & \GaussVIElementLeaders & \GaussVIAngle{34}{}{37}{}{28}{.35}\\
$F$ & \GaussVIElementLeaders & $0.2447424$\\
$G$ & \GaussVIElementLeaders & $0.4422071$
\end{tabular}
\end{center}
\endgroup
Die folgende Übersicht zeigt die Fehler dieser Elemente:
\begin{center}\small
\setlength{\tabcolsep}{4pt}
\newcommand{\GaussVIError}[2]{\makebox[11pt][c]{$#1$}\makebox[39pt][r]{$#2$}}
\begin{tabular}{c|c|c}
Opposition & \multicolumn{2}{c}{Unterschied}\\
von & der heliozentr. Länge & der geozentr. Breite\\\hline
1803 & \GaussVIError{-}{111''00} & \GaussVIError{-}{8''31}\\
1804 & \GaussVIError{+}{59.18} & \GaussVIError{-}{36.67}\\
1805 & \GaussVIError{+}{19.92} & \GaussVIError{+}{0.07}\\
1807 & \GaussVIError{+}{85.77} & \GaussVIError{+}{25.01}\\
1808 & \GaussVIError{+}{135.88} & \GaussVIError{+}{28.72}\\
1809 & \GaussVIError{-}{216.54} & \GaussVIError{+}{83.01}
\end{tabular}\end{center}
Den Hauptinhalt der vorliegenden Abhandlung bildet die Entwicklung der Methoden, mit denen diese verschiedenen Systeme von Elementen gefunden wurden. In dieser Hinsicht lässt sie sich als eine Art Ergänzung zu einigen Abschnitten der \textit{Theoria motus corporum coelestium} ansehen. Die Aufgabe, aus vier beobachteten Orten eines Himmelskörpers seine Bahn zu bestimmen, war in jenem Werk zwar schon ausführlich behandelt worden. Dort wurden zwei Lösungen angegeben, eine für den Fall, dass die Bahn noch ganz unbekannt ist, und eine für den Fall, dass die schon näherungsweise bekannte Bahn verbessert werden soll. Die zweite Lösung hätte also auch zur Bestimmung der drei obigen Systeme von Elementen verwendet werden können. Doch gerade im vorliegenden Fall, in dem die vier Beobachtungen Oppositionen sind, hielt es Professor \textsc{Gauß} für vorteilhafter, eine ganz andere Methode zu verwenden. Sie wird in der
\SourceLeaf{69}%
Abhandlung ausführlich erklärt und am Beispiel der tatsächlichen Berechnung des dritten Systems von Elementen weiter erläutert. Für eine nähere Beschreibung dieser Methode fehlt hier der Raum.

Das vierte System von Elementen wurde nach den Grundsätzen berechnet, die im 3. Abschnitt des 2. Buchs der \textit{Theoria motus corporum coelestium} entwickelt sind. Auch hierzu gibt die vorliegende Abhandlung mehrere Ergänzungen, die den Astronomen hoffentlich willkommen sein werden. Zunächst eine bequeme Berechnung der Differenzialänderungen der heliozentrischen Länge und der geozentrischen Breite aus den Differenzialänderungen der einzelnen Elemente. Sodann ein eigenes Verfahren, um die unbekannten Größen nach dem oben erwähnten Grundsatz zu bestimmen. Sind nämlich $w,w',w''$ usw. die vorgegebenen linearen Funktionen der unbekannten Größen $p,q,r$ usw., und soll die Summe $ww+w'w'+w''w''+$ usw. ein Minimum annehmen, so erhält man leicht ebenso viele lineare Gleichungen, wie es unbekannte Größen gibt. Aus ihnen sind die Unbekannten durch Elimination zu bestimmen. Bei einer einigermaßen großen Zahl von Unbekannten ist diese Elimination jedoch äußerst mühsam, da jede Gleichung alle Unbekannten enthält. Professor \textsc{Gauß} hat diese Arbeit erheblich abgekürzt. Zwar führt auch er die Lösung auf ebenso viele lineare Gleichungen zurück, wie es unbekannte Größen gibt. Seine Gleichungen sind jedoch so beschaffen, dass nur die erste alle Unbekannten enthält. Die zweite ist frei von $p$, die dritte von $p$ und $q$, die vierte von $p,q$ und $r$ usw. Daher macht es nur wenig Mühe, die Unbekannten in umgekehrter Reihenfolge zu bestimmen. Das Verfahren hat außerdem den Vorteil, dass man den kleinsten Wert von $ww+w'w'+w''w''+$ usw. im Voraus angeben und ihn zur Kontrolle der Rechnung mit dem Wert vergleichen kann, den man anschließend erhält, indem man die gefundenen Werte der Unbekannten in $w,w',w''$ usw. einsetzt.
\par\bigskip\centerline{\rule{.24\textwidth}{.3pt}}
